% Folder 144, batch 03 (archivists' pages 41-60).
% Pass: Opus 5.5 (claude-opus-5-5), 2026-10-03 — first pass, unchecked against the pages by a human.
%
% Skipped: 43, 45, 47 (administrative forms and letters, no mathematics), 48 (a cover in his hand,
% « 2-cartes cellulaires », used as the section heading).
\documentclass[11pt,a4paper]{article}
\input{../preamble/grothendieck}

\folder{144}
\batch{3}
\pages{41}{60}
\foldertitle{[Suite autour de Teichmüller dont] Teichmülleries, 1981-1982 : notes manuscrites (1981-1983, s.d.), lettres (1981, s.d.).}
\dating{1981-1983}
\watermark{Édition de démonstration}

\begin{document}

\page{41}
\note{p.~1 de l'auteur.}

\begin{tikzcd}[column sep=large]
  PM_{1,1} \arrow[d] & P\mathcal{D} \arrow[l] \arrow[d, "{(z_{1},z_{2}) \mapsto z_{2}/z_{1}}"] \\
  M_{1,1} & \mathcal{D} \arrow[l, "{\text{gal.\ groupe } \mathrm{Sl}(2,\mathbb{Z})}"']
\end{tikzcd}

\note{entre les deux flèches verticales, un mot abrégé lu \uncertain{cart}.}
\[
  P\mathcal{D} \simeq \{z_{1}, z_{2} \in \mathbb{C} \mid z_{2}/z_{1} \in \mathcal{D}\} \simeq \mathrm{Hom}^{+}(\mathbb{Z}^{2}, \mathbb{C})
\]
\note{au-dessous, une formule biffée et surchargée, où se lisent $\mathcal{D} \times \mathbb{C}^{*}$ ; l'exposant de $\mathrm{Hom}$ est lui-même récrit.}
\[
  \mathcal{D} \simeq \widetilde{M}_{1,1} \qquad \mathcal{D} = \{z \in \mathbb{C} \mid \Im z > 0\}
\]
(rigidification de Torelli \uncertain{stricte}, i.e.\ avec $\varphi(e_{1}, e_{2}) = +1$, $e_{2}/e_{1} = \tau \in \mathcal{D}$)

\drawing{Deux vecteurs issus d'un même point : $z_{1}$ (surchargé) vers le haut, $z_{2}$ vers la droite, avec une flèche courbe indiquant le sens de rotation.}

\note{avant $\mathbb{Z}^{2}$, un symbole biffé illisible.}
\[
  \mathbb{Z}^{2} \to \qquad
  \begin{pmatrix} a & b \\ c & d \end{pmatrix} \begin{pmatrix} x \\ y \end{pmatrix} = \begin{pmatrix} ax + by \\ cx + dy \end{pmatrix}
\]
\[
  \begin{pmatrix} a & b \\ c & d \end{pmatrix}(z) = \text{\struck{$\frac{dz + c}{bz + a}$}}\ \frac{az + b}{cz + d}
\]

$\mathbb{C}^{*} \times \mathrm{Sl}(2, \mathbb{Z})$ opère sur $P\mathcal{D}$ par $(z_{1}, z_{2}) \mapsto \lambda(az_{1} + bz_{2}, cz_{1} + dz_{2})$, \ill{} par passage au quotient, $\mathrm{Sl}(2, \mathbb{Z})$ opère sur $P\mathcal{D}/\mathbb{C}^{*} \simeq \mathcal{D}$, c'est induit par l'opération naturelle de $\mathrm{Gl}(2, \mathbb{R})$ sur $\hat{\mathbb{C}} = \mathbb{P}^{1}(\mathbb{C})$
\[
  \begin{pmatrix} a & b \\ c & d \end{pmatrix}(z) = \frac{az + b}{cz + d}
\]
$\mathrm{Sl}(2, \mathbb{R})$ invariant $\mathcal{D}$ (\uncertain{et même} $\mathrm{Sl}(2, \mathbb{R}) \simeq \mathrm{Aut}_{\mathbb{C}}(\mathcal{D})$), $\mathrm{Sl}(2, \mathbb{R}) \supset \mathrm{Sl}(2, \mathbb{Z})$.

\note{encadré par l'auteur :}
\begin{align*}
  M_{1,1}^{\mathrm{an}} &\simeq (\mathcal{D}, \mathrm{Sl}(2, \mathbb{Z})) \\
  PM_{11}^{\mathrm{an}} &\simeq (P\mathcal{D}, \mathrm{Sl}(2, \mathbb{Z}))
\end{align*}
\note{après $M_{1,1}^{\mathrm{an}}$ et $PM_{11}^{\mathrm{an}}$, un indice biffé.}
\begin{align*}
  M_{11}^{\mathrm{an}} &\simeq (\mathcal{D}, \mathrm{Gl}(2, \mathbb{Z})) \\
  PM_{11}^{\mathrm{an}} &\simeq (P\mathcal{D}, \mathrm{Gl}(2, \mathbb{Z}))
\end{align*}
\begin{align*}
  T_{11} &\simeq \mathrm{Sl}(2, \mathbb{Z}) \\
  ST_{11} &\simeq \widetilde{\mathrm{Sl}}(2, \mathbb{Z})
\end{align*}

$\mathrm{Gl}^{\pm 1}(2, \mathbb{R})$ opère sur $P\mathcal{D}$ par
\[
  (z_{1}, z_{2}) \mapsto (a\tau^{\alpha}z_{1} + b\tau^{\alpha}z_{2},\ c\tau^{\alpha}z_{1} + d\tau^{\alpha}z_{2})
\]
\[
  u_{g} = \tau^{\alpha(g)} \cdot g \qquad
  \begin{cases}
    \alpha(g) \in \mathbb{Z}/2\mathbb{Z} \\
    \alpha(g) = 0 & \text{si } \det g = 1 \\
    \phantom{\alpha(g)} = 1 & \text{si } \det g = -1
  \end{cases}
\]
\note{ici $\tau$ désigne sans doute la conjugaison complexe, non le $\tau \in \mathcal{D}$ plus haut ; les exposants $\alpha$ sont surchargés.}

\page{42}
\note{p.~2 de l'auteur.}
Section $\tau \mapsto (\tau, 1)$ de $P\mathcal{D}$ sur $\mathcal{D}$, d'où \uncertain{iso}
\note{au-dessus, une formule biffée commençant par $P\mathcal{D}$.}
\begin{align*}
  \mathcal{D} \times \mathbb{C}^{*} &\to P\mathcal{D} \\
  (\tau, \lambda) &\mapsto \lambda(\tau, 1) = (\lambda\tau, \lambda) \overset{\text{déf}}{=} [\tau, \lambda]
\end{align*}
on trouve
\[
  \begin{cases}
    \lambda'[\tau, \lambda] = [\tau, \lambda'\lambda] \\
    \begin{pmatrix} a & b \\ c & d \end{pmatrix}([\tau, \lambda]) = \left[\dfrac{a\tau + b}{c\tau + d},\ \lambda(c\tau + d)\right]
  \end{cases}
\]
\begin{align*}
  \widetilde{PM}_{11}^{\mathrm{an}} \simeq \mathcal{D} \times \mathbb{C} &\to P\mathcal{D} \\
  (\tau, z) &\mapsto (\exp 2\pi i z)(\tau, 1) = [\tau, \exp 2i\pi z]
\end{align*}
\[
  \overset{\omega^{2}}{\mathbb{Z}} \subset \widetilde{\mathrm{Sl}}(2, \mathbb{Z})
\]
\note{au-dessus, un mot biffé.}
opère sur $\mathcal{D} \times \mathbb{C}$ de façon compatible \uncertain{avec} l'action sur $P\mathcal{D}$ et $\omega^{2} \in \mathbb{Z}$ opérant par
\[
  (\tau, z) \mapsto (\tau, z + 1)
\]

Pts fixes de $z \mapsto \dfrac{az + b}{cz + d}$ donnés par
\[
  cz^{2} + (d - a)z - b = 0 \qquad \Delta = (d - a)^{2} + 4bc = (a + d)^{2} - 4 \quad \text{car } bc = ad - 1
\]
Ici $a, b, c, d \in \mathbb{Z} \subset \mathbb{R}$, \uncertain{on a des} \struck{\ill{}} \uncertain{conditions} \ill{} $\exists$ $z$ pt fixe $\in \mathcal{D}$ i.e.\ $z \notin \mathbb{R}$ i.e.\ on \ill{} \ill{} $\Delta < 0$ i.e.\ $|a + d| \leq 1$
\note{un morceau d'étiquette imprimée (« GP », « THEMATIQUE ») est collé sur le bas de la page et cache une partie de la ligne.}

Deux cas \struck{$a + d$}
1°) $a + d = 0$, $d = -a$, $d - a = -2a$
\[
  z \mapsto \frac{az + b}{cz - a} \qquad -a^{2} - bc = 1, \quad bc = -1 - a^{2}
\]
\[
  cz^{2} - 2az - b = 0 \qquad z = \frac{a \pm 2i}{2c} \qquad \text{On trouve (4)}
\]
pour $a = 0$, \struck{\ill{}} $d = 0$, $b = -1$, $c = 1$, $z = i$
\note{$z = \frac{a \pm 2i}{2c}$ est lu tel quel ; le discriminant donne $z = \frac{a \pm i}{c}$.}

\page{44}
\note{p.~3 de l'auteur.}
2°) $a + d = 1$ i.e.\ $d = 1 - a$, par ex.\ $a = 0$, $d = 1$, $bc = -1$, $b = 1$, $c = -1$
\[
  \rho = \begin{pmatrix} 0 & 1 \\ -1 & 1 \end{pmatrix}, \qquad z = \frac{1}{2} + i\frac{\sqrt{3}}{2} = \exp\frac{2i\pi}{6} = -j
\]
\note{le dernier membre, lu $-j$, est douteux.}

\page{46}
\note{p.~4 de l'auteur.}
\drawing{Le demi-plan supérieur : axe réel, axe imaginaire, demi-cercle unité ; la bande verticale hachurée au-dessus de l'arc qui va de $i$ au point marqué $j$ (domaine fondamental).}
\[
  \rho = \begin{pmatrix} 0 & 1 \\ -1 & 1 \end{pmatrix} \qquad \rho^{3} = \omega = \begin{pmatrix} -1 & 0 \\ 0 & -1 \end{pmatrix}
\]
\[
  \sigma = \begin{pmatrix} 0 & +1 \\ -1 & 0 \end{pmatrix} \qquad \sigma^{2} = \omega, \quad \sigma^{-1} = \omega\sigma
\]

On a choisi les signes de telle façon que a) $\rho^{3} = -1$ i.e.\ $\rho$ d'ordre $6$ i.e.\ il engendre le groupe d'isotropie \add{de $j$} (mais cela laisse le choix entre $\rho$ et $\rho^{-1}$) b) $\sigma$ soit la rotation d'angle $-\frac{\pi}{2}$ dans $\mathbb{R}^{2}$ (pour la structure euclidienne habituelle) et que $\rho$ soit une rotation d'angle $-2\pi/6$ (pour une métrique invariante convenable) --- la chose importante dans cette convention, c'est que pour $\rho$ et $\sigma$, le sens de rotation (ici, négatif) soit le même. Quand on interprète \struck{par} les matrices $\begin{pmatrix} a & b \\ c & d \end{pmatrix} \in \mathrm{Sl}(2, \mathbb{Z})$ comme opérant sur des \struck{\ill{}} \add{couples} $(z_{1}, z_{2})$ de vecteurs de $\mathbb{C}$, engendrant un sous-groupe discret de $\mathbb{C}$ (avec $z_{1}/z_{2} \in \mathcal{D}$), alors l'action de $\rho$ et de $\sigma$ sur $z_{1}, z_{2}$ est \og positive \fg, en ce sens que $(z_{1}, u(z_{1}))$ et $(z_{2}, u(z_{2}))$ forment des $\mathbb{R}$-bases \emph{directes} de $\mathbb{C} \simeq \mathbb{R}^{2}$.

Voyons le groupe des relèvements de $\mathrm{Sl}(2, \mathbb{Z})$ opérant sur $P\mathcal{D}$ sur $\widetilde{P\mathcal{D}}$, qui est donc une extension centrale de $T_{11} = \mathrm{Sl}(2, \mathbb{Z})$ par $\mathbb{Z}$, \struck{\ill{}} \add{\uncertain{\ill{}}} \add{\uncertain{extension}} qu'on va écrire $ST_{11}$ \struck{$\mathrm{Sl}(2, \mathbb{Z})$}. Je désignerai par \struck{\ill{}} $\varpi$ l'élément de \struck{$\mathrm{Sl}(2, \mathbb{Z})$} $ST_{11}$ \uncertain{défini} par le générateur de $\mathbb{Z}$, opérant sur $P\mathcal{D}$ par $(\tau, z) \mapsto (\tau, z + 1)$
\note{la phrase reste en suspens au bas de la page ; la suite n'est pas dans ce lot.}

\section{\uncertain{Rétrospection} sur les cartes cellulaires}
\note{la page 48, une chemise par ailleurs vierge, porte de sa main, au crayon : « 2-cartes cellulaires ». Le titre de section est celui qu'il écrit en tête de la page 49.}

\page{49}
\[
  (X, G) = \Pi_{1}(X, G)
\]
\[
  \mathop{\mathrm{Hom}}_{(X,G)}(x, y) = \{(\ell, g) \mid \ell : gx \to y,\ g \in G\}
\]
\note{sous la flèche $gx \to y$ : $\in \mathrm{FC}\,\Pi_{1}X$ (lecture douteuse).}

\begin{tikzcd}[column sep=small, row sep=small]
  gx \arrow[r, "\ell"] & y & \\
  & g'y \arrow[r, "\ell'"] & z
\end{tikzcd}

\note{encadré par l'auteur :}
\[
  (\ell', g')(\ell, g) = (\ell' \circ g'(\ell),\ g'g)
\]

\page{50}
\note{table des matières de la \og Rétrospection \fg, avec les renvois aux pages de l'auteur.}
\begin{enumerate}
  \item Description combinatoire des cartes, via $\mathfrak{G}_{2}$ et $\mathfrak{G}_{2}^{0}$ --- p.~1
  \item Cartes triangulaires pondérées, et groupes $\pi_{03}$, $\tilde\pi_{03}$ --- p.~4
  \item Relation avec revêtements de $U_{0,3}$ et courbes algébriques --- p.~5'
  \item Conditions $\rho_{i}^{\nu_{i}} = 1$ et groupes fondamentaux à signature --- p.~9'
  \item Sur l'existence d'une pondération d'une carte triangulaire --- p.~12'
  \item Carte associée à un polygone convexe, et groupes fondamentaux associés --- p.~15
  \item Relation aux cartes $n$-polygonales pondérées --- p.~22
  \item Introduction de $P\Sigma_{p_{n}}^{\times}$ et du \uncertain{$2$-groupoïde} fondamental --- p.~29
\end{enumerate}

\page{51}
\note{p.~1 de l'auteur ; c'est le début de la \og Rétrospection \fg{} annoncée par la table de la page 50.}
\uncertain{Rétrospection} sur les cartes (finies)

(1) Définition initiale \add{d'une \emph{carte} (topologique) cellulaire} : surface compacte (pas nécessairement orientable) $X$, avec \struck{\ill{}} fermée $K$, loc.\ \struck{homéom.}\ à une étoile (\og 1-complexe topologique \fg), plus hypothèses \og cellulaires \fg : a) $K$ n'a pas de composantes connexes homéom.\ à $S^{1}$ \struck{ou vide} \add{ni de point isolé} b) les comp.\ connexes de $X \setminus K$ sont des disques ouverts i.e.\ homéom.\ à $\mathring{\mathbb{D}} = \{z \in \mathbb{C} \mid |z| < 1\}$. On en fait une \og catégorie isotopique \fg, et à l'aide de la notion de biarc (p.ex.), on établit une équivalence de catégories
\[
  (1)\qquad \text{Cartes isot.} \xrightarrow{\ \approx\ } \mathfrak{G}_{2}\text{-ensembles finis spéciaux}
\]
\note{devant « Cartes isot. », un premier départ biffé (\struck{Catisot}) ; « finis » est ajouté par lui au-dessus de la ligne, et « spéciaux » est entre guillemets.}
où
\[
  (2)\qquad \mathfrak{G}_{2} \simeq \{\tau_{0}, \tau_{1}, \tau_{2} \mid \tau_{0}^{2} = \tau_{1}^{2} = \tau_{2}^{2} = (\tau_{0}\tau_{2})^{2} = 1\}
\]
\note{sous $\tau_{2}$, de sa main : « ou $\tau_{\infty}$ ».}
(NB Par la suite, on notera $\tau_{\infty}$ au lieu de $\tau_{2}$, notation \struck{\ill{}} mieux appropriée pour les relations avec la Géom.\ Alg. La notation $\tau_{2}$ est plus appropriée pour les généralisations en dim.\ supérieure …), et où les $\mathfrak{G}$-ens \og spéciaux \fg{} sont ceux où les involutions
\[
  \tau_{0}, \tau_{1}, \tau_{2} \quad \text{et} \quad \sigma = \tau_{0}\tau_{2}
\]
opèrent sans points fixes. (NB.\ Par la suite, on posera $\rho_{1} = \tau_{0}\tau_{\infty}$, en même temps qu'on écrit $\tau_{\infty}$ pour $\tau_{2}$)

On introduit
\[
  (3)\qquad \mathfrak{G}_{2}^{0} = \mathrm{Ker}\,(\mathfrak{G}_{2} \to \mathbb{Z}/2\mathbb{Z}), \quad \tau_{i} \mapsto 1 \bmod 2
\]
et on a
\[
  (4)\qquad \mathfrak{G}_{2}^{0} = \{\rho_{0}, \rho_{1}, \rho_{\infty} \mid \rho_{\infty}\rho_{1}\rho_{0} = \rho_{1}^{2} = 1\}
\]
\note{entre les deux signes $=$, un symbole biffé.}

\marginal{b) i.e.\ que $K$ \ill{} \ill{} \ill{} \ill{} \ill{} il y a \ill{} $x_{0}$ \ill{} $X$ \ill{} \ill{} \ill{} \ill{} \ill{} \ill{}}

\marginal{La notation $\sigma_{2}$ est plus commode \ill{} \ill{} \ill{} \ill{} \ill{} $\sigma_{i}$ \ill{} \ill{} \ill{} \ill{} \ill{} \ill{} \ill{}}
\note{deux notes marginales écrites en diagonale le long du bord gauche ; seuls quelques mots se laissent lire.}

\page{52}
avec
\[
  (5)\qquad \rho_{\infty} = \tau_{1}\tau_{0}, \quad \rho_{1} = \tau_{0}\tau_{\infty}, \quad \rho_{0} = \tau_{\infty}\tau_{1}
\]

La relation avec les \og repères \fg, \uncertain{en tant que} simplexes topologiques de la subdivision barycentrique, est la suivante :

\drawing{Un morceau de carte trivalente (traits gras) et, en traits fins au crayon, sa subdivision barycentrique. Au sommet $s$ de la carte, le long de l'arête $a$ (flèche $\to a$) et dans la face $f$, le triangle hachuré de sommets $s$, $\tilde a$ (centre de $a$) et $\tilde f$ (centre de $f$) est le repère $r$ ; une flèche courbe indique le sens de rotation du repère, une flèche verticale l'orientation normale.}

Le biarc \add{associé au repère $r$} donne \uncertain{comme} orientations tangentielles celle qui \emph{sort} des sommets, comme orientation normale celle qui \emph{entre} dans la face $f$ (\og on sort des sommets pour entrer dans la face $f$, en passant par l'arête $a$ \fg). Les trois sommets de $r$ sont le sommet $s = \tilde s$, le \og centre \fg{} $\tilde a$ de $a$, et le \og centre \fg{} $\tilde f$ de $f$, l'orientation de $r$ est celle qui tourne autour de $r$ dans le sens $s\,\tilde a\,\tilde f$ (comme le dit la \uncertain{localisation imagée} plus haut). C'est aussi celle pour laquelle l'\struck{\ill{}} orientation \add{de l'arête $a$} \ill{} \ill{} autour des sommets, \ill{} orientations tangentielles \ill{}, indique le sens de rotation \og \ill{} \fg{} de $a$ pour l'orientation \ill{} et pour laquelle l'orientation tangentielle indique le sens de rotation autour de $f$ (\og face \uncertain{extrémité} \fg{} de $a$ pour l'orientation normale du biarc). Les sens de rotation \og \uncertain{dérivés} \fg{} (par l'orientation définie par $r$) autour des trois sommets de $r$ sont indiqués sur la figure. Sur les \uncertain{positions formées par} $r$, \ill{} les trois sommets, le \struck{triangle} \add{\uncertain{simplexe}} topologique $r$, \ill{} sur la figure le sens de l'orientation \ill{} des arêtes \add{$\tilde s\tilde a$, $\tilde a\tilde f$, $\tilde f\tilde s$} \ill{} à leur sens, \ill{} des positions le long des \ill{} qui tournent autour des sommets.

\page{53}
\note{p.~2 de l'auteur.}
La \struck{complexe} \add{carte} \emph{duale} d'une carte (\struck{\ill{}} \uncertain{combinatoirement}) a le même \uncertain{ensemble} de repères \og triangulés \fg{} \struck{par} les mêmes repères, mais le rôle des sommets de type $0$ et $\infty$ échangés : les sommets de $C$ deviennent \add{et inversement} \struck{les \ill{}} les faces de $\check C$, les faces de $C$ correspondent aux sommets de $\check C$, enfin les arêtes de $C$ et celles de $\check C$ se correspondent (avec relation d'incidence). On notera que le repère commun $r$, \struck{\ill{}} quand il est regardé comme repère pour $\check C$ \add{au lieu de} \struck{\ill{}} $C$, reçoit des orientations \og \emph{opposées} \fg, tant \uncertain{dans} les orientations \uncertain{opposées} sur les arêtes et pour les sens de rotation autour des sommets.

Une carte $X$ est \emph{orientée} si la surface sous-jacente est orientée, ce qui \uncertain{revient} à ce que, \uncertain{sur} l'ensemble $\mathcal{R}(X)$ des repères de $X$, on se donne une subdivision en deux parties complémentaires
\[
  (5)\qquad \mathcal{R}(X) = \mathcal{R}^{+}(X) \amalg \mathcal{R}^{-}(X)
\]
telles que l'on ait
\[
  (6)\qquad \sigma_{i}(\mathcal{R}^{+}(X)) \subset \mathcal{R}^{-}(X) \quad \text{pour } i \in \{0, 1, \infty\}
\]
\uncertain{(6 bis)} \ill{} \uncertain{égalité}, ce qui implique
\begin{align*}
  (7)&\qquad \sigma_{i}(\mathcal{R}^{+}(X)) = \mathcal{R}^{+}(X), \quad \sigma_{i}(\mathcal{R}^{-}(X)) = \mathcal{R}^{-}(X) \quad i \in \{0, 1, \infty\} \\
  (7\text{ bis})&\qquad \rho_{i}(\mathcal{R}^{+}(X)) = \mathcal{R}^{+}(X), \quad \rho_{i}(\mathcal{R}^{-}(X)) = \mathcal{R}^{-}(X) \quad i \in \{0, 1, \infty\}
\end{align*}
\note{la ligne (7) est lue telle qu'elle paraît, mais elle contredit (6), dont elle devrait être la forme forte ($\sigma_{i}$ échangeant $\mathcal{R}^{+}$ et $\mathcal{R}^{-}$) ; les exposants de la première ligne sont douteux. Sur ces pages les involutions sont écrites tantôt $\tau_{i}$, tantôt $\sigma_{i}$ ; la transcription suit la page.}

Alors aussi $\mathcal{R}^{+}(X)$ est un $\mathfrak{G}_{2}^{+}$-ensemble : \ill{}, et $X \mapsto \mathcal{R}^{+}(X)$ donne une équivalence de catégories
\[
  (8)\qquad \text{Cartes isot.\ or.} \xrightarrow{\ \approx\ } \mathfrak{G}_{2}^{0}\text{-ens.\ spéciaux finis}
\]
\note{« spéciaux » : lecture douteuse.}

\marginal{NB On utilise une notion de \og morphisme \fg{} de cartes qui n'est pas tellement triviale (\ill{} \ill{} \uncertain{pour les} \ill{}).}

\page{54}
\drawing{Une figure à l'encre et au crayon : le complexe trivalent $K$ en traits gras (sommets marqués $0$), sa subdivision barycentrique en traits fins (sommets marqués $0$, $1$, $\infty$), et dix triangles-repères étiquetés $r$ (hachuré), $\tau_{0}r$, $\tau_{1}r$, $\tau_{\infty}r$, $\tau_{0}\tau_{1}r = \rho_{\infty}r$, $\tau_{1}\tau_{0}r = \rho_{\infty}r$, $\tau_{\infty}\tau_{1}r = \rho_{0}r$, $\tau_{0}\tau_{\infty}r = \rho_{0}^{-1}r$, $\tau_{0}\tau_{\infty}r = \tau_{\infty}\tau_{0}r = \rho_{1}r$, avec les sens de rotation et les flèches d'orientation des arêtes.}

\marginal{le complexe topologique $K$ est un \uncertain{graphe tracé} en traits gras}

\marginal{NB $\sigma_{i}$ change le sommet de type $i$ de $r$, et fixe les deux autres sommets}

Les opérations $\tau_{0}, \tau_{1}, \tau_{\infty}$ sont marquées sur la figure, ce sont les \og symétries \fg{} p.r.\ aux côtés du triangle $r$ opposés aux sommets $s = S^{0}_{r}$, $\tilde a = S^{1}_{r}$, $\tilde f = S^{\infty}_{r}$ de $r$ (sommets dits de type $0$ (ou \og type-sommet \fg), type $1$ (ou \og type arête \fg), type $\infty$ (ou $2$) (ou \og type-face \fg)). On notera que les orientations \uncertain{induites} par deux repères \og adjacents \fg{} le long de l'arête commune (ce qui \uncertain{capture} \uncertain{par les} \add{deux} orientations de la surface \uncertain{ambiante}, au \uncertain{voisinage} de celle-ci, définies par $r$ et $r'$ respectivement), sont toujours dans les sens \og \uncertain{opposés} \fg. Par contre, les \add{sens de rotation} \struck{orientations} autour des sommets communs définis \struck{\ill{}} respectivement par deux repères adjacents, sont \emph{opposés}. La figure montre \uncertain{dix} repères, savoir $r$ et ceux qui s'en déduisent \struck{\ill{}} les opérations $\tau_{i}$ ($i \in \{0, 1, \infty\}$), et en \add{recomposant} \struck{\ill{}} une fois
\begin{align*}
  & r \qquad \tau_{0}r, \tau_{1}r, \tau_{\infty}r \qquad \tau_{1}\tau_{0}r = \rho_{\infty}r, \ \tau_{0}\tau_{1}r = \rho_{\infty}^{-1}r \\
  & \phantom{r \qquad \tau_{0}r, \tau_{1}r, \tau_{\infty}r \qquad} \tau_{0}\tau_{\infty}r = \rho_{1}r = \tau_{\infty}\tau_{0}r = \rho_{1}^{-1}r \\
  & \phantom{r \qquad \tau_{0}r, \tau_{1}r, \tau_{\infty}r \qquad} \tau_{\infty}\tau_{1}r = \rho_{0}r, = \tau_{1}\tau_{\infty}r = \rho_{0}^{-1}r
\end{align*}
\note{les signes $=$ entre $\rho_{1}r$ et $\tau_{\infty}\tau_{0}r$, et entre $\rho_{0}r$ et $\tau_{1}\tau_{\infty}r$, semblent récrits ou barrés à l'encre : la page n'affirme sans doute pas ces égalités ; les indices des troisièmes et quatrièmes termes sont eux aussi surchargés.}

(il y aurait $10$ repères, au lieu de $9 = 1 + 3 + (2 + 1 + 2)$, s'il n'y avait ici la relation
\[
  \tau_{0}\tau_{\infty} = \tau_{\infty}\tau_{0} \quad \text{i.e.} \quad \rho_{1} = \rho_{1}^{-1} \quad \text{i.e.} \quad \rho_{1}^{2} = 1 .
\]

\page{55}
\note{p.~3 de l'auteur.}
\ill{} maintenant \og spéciaux \fg{} signifie que $\rho_{1} = \tau_{0}\tau_{\infty}$ opère sans points fixes.

Si on a une carte orientée, on \struck{considère} la carte \emph{duale} munie de la \emph{même orientation de la surface sous-jacente commune}, est appelée la carte \emph{orientée duale}. On a alors
\[
  (9)\qquad \mathcal{R}^{+}(\check C) = \mathcal{R}^{-}(C), \quad \mathcal{R}^{-}(\check C) = \mathcal{R}^{+}(C)
\]
i.e.\ les repères directs pour l'une des deux cartes sont ceux qui sont inverses pour l'autre.

On se libère de la condition de restriction aux $\mathfrak{G}_{2}^{0}$-ensembles finis \emph{spéciaux}, en introduisant une notion un peu plus générale de carte cellulaire, \struck{en \uncertain{admettant}} \add{en introduisant} la notion \struck{\ill{}} d'\emph{arête pliée} --- Techniquement, \ill{}, on se donne, \uncertain{outre} $K$, une partie finie $V$ (ens.\ des sommets) $S \subset K$, qui contient tous les sommets topologiques, sauf éventuellement certains sommets d'ordre $1$ ou points-bords de $K$.

\drawing{Un graphe (arbre) en traits gras épaissis, entouré d'un contour fin qui en dessine le voisinage (le disque $D(X/K)$) ; sommets marqués en gras, quelques points sur le bord. Une flèche désigne une extrémité : \og arête repliée \fg, avec le petit croquis d'un \og modèle intuitif \fg{} (un Y épaissi dont une branche est repliée sur elle-même).}

\marginal{$\leftarrow$ arête repliée (sur le modèle intuitif ci-contre)}

\marginal{les sommets marqués sont marqués en gras, le contour du disque $D(X/K)$ en trait fin, les sommets sur le bord de \uncertain{le} disque\note{sic} \struck{sont traits} figurés \ill{}}

On reformule directement les axiomes en termes de $S \subset K \subset X$. Du point de vue des repères le long d'une arête pliée il y a deux repères (au lieu de quatre dans le cas précédemment étudié)

\page{56}
\drawing{Un arbre trivalent dont les sommets portent les types $0$ et $1$ ; à droite une arête repliée qui aboutit en un sommet de type $\infty$, avec les deux repères $r$ et $r'$ de part et d'autre (sens de rotation opposés) dans un contour fin ; en bas une région hachurée autour d'un sommet $\infty$ où figurent $\sigma_{\infty}r_{1}$, $\sigma_{0}r_{1}$, $\sigma_{1}r_{1}$ et $r_{1}$.}

\[
  (10)\qquad
  \begin{cases}
    r' = \tau_{\infty}r = \tau_{0}r \\
    \text{i.e.}\quad \tau_{0}\tau_{\infty}r = r \quad \text{i.e.}\quad \rho_{1}r = r
  \end{cases}
\]

Les arêtes \uncertain{repliées} correspondent $1$-$1$ aux \og sommets de type $1$ \fg{} (de la subdivision barycentrique) qui sont topologiques (\uncertain{univ.} d'ordre $1$) de $K$. Et elles correspondent $1$-$1$ aux points fixes de $\rho_{1} = \sigma_{0}\sigma_{\infty}$ opérant sur $\mathcal{R}$.

Pour un pt.\ bord de $K$ qui est un sommet marqué de $K$ (i.e.\ un sommet de type $1$ de la subdivision barycentrique) \ill{} la situation ci-contre

\drawing{Un sommet marqué $0$ au bord de $K$, d'où partent des arêtes vers des sommets $0$ et $\infty$ ; trois boucles fines emboîtées autour de l'arête qui aboutit en ce sommet, étiquetées $\tau_{0}r$, $r$ et $\tau_{\infty}\tau_{0}r = \tau_{1}\tau_{0}r$ (lecture douteuse), avec flèches de rotation.}

, il y a bien \emph{4} repères distincts adjacents à l'arête qui aboutit en ce sommet --- si on prend l'un des deux qui \struck{\uncertain{combinent}} \add{\ill{}} \og \add{pli} \fg{} \struck{\og pli \fg} \struck{sommet} \ill{}, on aura
\[
  (11)\qquad \tau_{1}r = \tau_{\infty}r \quad \text{i.e.} \quad \rho_{0}r = r
\]
ce qui caractérise les sommets \uncertain{envisagés} (\ill{} \ill{}), en corr.\ biunivoque avec les paires de repères $\{r, r'\}$ telles que $r' = \tau_{1}r = \tau_{\infty}r$, d'où $r = \tau_{1}r' = \tau_{\infty}r'$. \struck{Condition} La condition (11) n'impliquerait pas que

\page{57}
\note{p.~4 de l'auteur. La phrase continue celle de la page précédente.}
puisse avoir \struck{\ill{}} lieu
\[
  \tau_{1}r = \tau_{\infty}r = \tau_{0}r
\]
i.e.\ qu'on soit dans le cas (10) (\uncertain{qui me semble} \ill{} : \uncertain{qu.} les $\sigma_{i}$ opèrent sans pts fixes) cela signifie que le point de \emph{type 1} du repère \ill{} est un pli (le point de type $0$ ne pourrait l'être, par définition).

Grâce à l'introduction des arêtes pliées, dans \uncertain{la} équivalence\note{sic} de catégories (8), il n'est plus nécessaire d'introduire l'hypothèse que $\rho_{1}$ opère sans pt fixe --- et dans l'équivalence (1), il n'y a plus lieu de supposer que $\rho_{1}$ (disons $\sigma$) opère sans pt fixe --- par contre la condition que les $\sigma_{i}$ opèrent sans points fixes demeure provisoirement (mais il serait facile maintenant de s'en affranchir aussi cf plus bas …)).

(2) On a dans ces considérations un rôle particulier pour l'indice $1 \in \{0, 1, \infty\}$, du fait que l'on a
\[
  \rho_{1}^{2} = 1 ,
\]
on peut s'en défaire, en introduisant
\[
  (12)\qquad \tilde\pi_{03} = \{\tau_{0}, \tau_{1}, \tau_{\infty} \mid \tau_{0}^{2} = \tau_{1}^{2} = \tau_{\infty}^{2} = 1\}
\]
(sans condition $(\tau_{0}\tau_{\infty})^{2} = 1$ !), et le sous-groupe d'indice $2$
\[
  (13)\qquad \pi_{03} = \{\rho_{0}, \rho_{1}, \rho_{\infty} \mid \rho_{\infty}\rho_{1}\rho_{0} = 1\}
\]

\marginal{$\pi_{03} = \mathrm{Ker}\,(\tilde\pi_{03} \to \mathbb{Z}/2\mathbb{Z})$, $\sigma_{i} \mapsto 1 \bmod 2$ (« Ker » : lecture douteuse)}

\page{58}
en prenant encore, comme dans (5)
\[
  (14)\qquad \rho_{\infty} = \tau_{1}\tau_{0}, \quad \rho_{1} = \tau_{0}\tau_{\infty}, \quad \rho_{0} = \tau_{\infty}\tau_{1} .
\]
On a donc
\[
  (15)\qquad \mathfrak{G}_{2} \simeq \tilde\pi_{03}/\langle\rho_{1}^{2}\rangle, \qquad \mathfrak{G}_{2}^{0} \simeq \pi_{03}/\langle\rho_{1}^{2}\rangle
\]

\struck{Cette fois} \add{Ici} \struck{\ill{}} $\tilde\pi_{03}$ apparaît comme le groupe \og classifiant \fg{} pour la structure de \og \emph{carte triangulaire} \emph{pondérée} \fg{} (par $\{0, 1, \infty\}$), i.e.\ une carte $\mathcal{C} = (X, K)$ \struck{l'ensemble} avec une application
\[
  (16)\qquad S \longrightarrow \{0, 1, \infty\}
\]
\note{sous $S$, en petit : \og ens.\ des sommets (ordinaires) de la carte (\uncertain{vue} d'une subd.\ barycentrique !) \fg.}
ayant la propriété que pour toute face $f$, l'application (16) induit une bijection de l'ensemble des sommets de $f$ avec $\{0, 1, \infty\}$. L'ensemble des \emph{repères} (au sens de la théorie \add{triangulaire} pondérée, au cas de possibilité \ill{} \ill{} \ill{} \og repères pondérés \fg) de $X$ est par définition l'ensemble de ses faces, et $\tau_{i}$ ($i \in \{0, 1, \infty\}$) est le \struck{\ill{}} repère symétrique p.r.\ au \struck{\ill{}} côté du repère opposé au sommet de type $i$. L'\emph{ordre} des sommets est pair \ill{}, \ill{} \emph{ordre réduit} d'un sommet le quotient de l'ordre par deux.

\page{59}
\note{p.~5 de l'auteur.}
\drawing{Une figure de dix triangles autour d'un triangle central $r$ (sommets étiquetés $0$, $1$, $\infty$ de manière à alterner), chacun marqué du repère qu'il représente et de son sens de rotation : $\tau_{0}r$, $\tau_{1}r$, $\tau_{\infty}r$, $\tau_{1}\tau_{\infty}r = \rho_{0}^{-1}r$, $\tau_{\infty}\tau_{1}r = \rho_{0}r$, $\tau_{0}\tau_{1}r = \rho_{\infty}^{-1}r$, $\tau_{1}\tau_{0}r = \rho_{\infty}r$, $\tau_{0}\tau_{\infty}r = \rho_{1}r$, $\tau_{\infty}\tau_{0}r = \rho_{1}^{-1}r$ ; les arêtes portent des flèches d'orientation.}

Les propriétés d'orientation de la figure sont les mêmes que celles explicitées p.~2', la figure est plus jolie ou du moins plus symétrique. Néanmoins, dans des cas particuliers certains des $10$ \struck{\ill{}} repères de la figure peuvent coïncider (ex : octaèdre, ou exemple plus dégénéré encore, le cas \og orienté universel \fg

\drawing{Une sphère dont l'équateur porte les trois points $0$, $1$, $\infty$.}

). Ici l'équivalence de catégories ad hoc
\[
  (17)\qquad \text{cartes isot.\ triangulaires pondérées} \xrightarrow{\ \approx\ } \tilde\pi_{03}\text{-ensembles finis spéciaux}
\]
\note{« spéciaux » : lecture douteuse.}
est en termes des ens.\ finis $\mathcal{R}$ : \ill{} involutifs $\tau_{i}$ ($i \in \{0, 1, \infty\}$) \emph{sans points fixes}, sans plus.

Définissons les cartes triangulaires pondérées orientées comme des cartes ordinaires, par une orientation de la surface top.\ sous-jacente, on voit que cette notion se définit encore combinatoirement par la donnée d'un sous-ens.\ $\mathcal{R}^{+}(X)$ de l'ens.\ $\mathcal{R}(\alpha)$ des repères pondérés de $X$, satisfaisant (6 bis), i.e.\ \add{(7 bis)} (7) \ill{} \ill{} (6). On a alors
\[
  \mathcal{R}^{-}(X) = \mathcal{R}(X) \setminus \mathcal{R}^{+}(X) .
\]

\page{60}
\[
  (18)\qquad \text{Cartes (isotop.) triangulaires pondérées orientées} \xrightarrow{\ \approx\ } (\pi_{03}\text{-Ens}) .
\]
\note{au-dessus de la flèche : $X \longmapsto \mathcal{R}^{+}(X)$.}

\marginal{NB \ill{} \ill{} \ill{} $\mathcal{R}$ \ill{} \ill{} $\pi_{03}$ \ill{}}

La notion de \og carte duale \fg{} se remplace ici par la notion de transformée d'une carte par une \og réindexation \fg{} des sommets, via un
\[
  g \in \mathfrak{S}_{3} \overset{\text{déf}}{=} \mathfrak{S}_{\{0, 1, \infty\}}
\]
--- \ill{} induits \ill{} \ill{}, \uncertain{cela est} trivial, \ill{} \ill{} \ill{} la formule
\[
  (18)\qquad \mathcal{R}^{+}(X^{g}) = \mathcal{R}^{\mathrm{sg}(g)}(X) \qquad g \in \mathfrak{S}_{3}
\]
\note{à droite : \og $\mathrm{sg}\,g$ la signature de $g$ \fg. Le numéro (18) est donné deux fois sur la page.}

(3) Par les considérations précédentes, et notamment par les isom.\ (15)
\begin{align*}
  \mathfrak{G}_{2} &\simeq \tilde\pi_{03}/\langle\rho_{1}^{2}\rangle \\
  \mathfrak{G}_{2}^{+} &\simeq \pi_{03}/\langle\rho_{1}^{2}\rangle ,
\end{align*}
la cat.\ \add{isotopique} des cartes finies (avec arêtes pliées) apparaît comme équivalente à une sous-catégorie pleine de la catégorie \add{isotopique} des cartes finies triangulaires pondérées, et \ill{} itou pour la version orientée. D'autre part, on a
\[
  (19)\qquad \pi_{03} \simeq \pi_{1}(U_{0,3}, P_{+}^{*})
\]
\note{l'exposant de $P_{+}$ est d'une lecture douteuse.}
où
\[
  (20)\qquad U_{0,3} = \mathbb{P}^{1}(\mathbb{C}) \setminus \{0, 1, \infty\}
\]
\struck{$P_{+}$ et} $P_{+}$ étant un pt \struck{quel} de l'hémisphère supérieur --- il est commode de prendre
\[
  P_{+} = -\zeta, \qquad \zeta = \exp 2i\pi/3 .
\]
\note{l'argument continue au-delà de la fin du lot.}

\end{document}
